\documentclass[11pt,a4paper]{article}
\usepackage[utf8]{inputenc}
\usepackage{amsmath}
\usepackage{amsfonts}
\usepackage{amssymb}
\usepackage{graphicx}
\usepackage{hyperref}
\usepackage{authblk}

\title{Differentiable Semi-Tensor Product Networks Reveal Holographic Vulnerabilities in Gene Regulatory Circuits}
\author[1]{Yaseen Khalil}
\affil[1]{AI Architect, Independent Researcher}
\date{\today}

\begin{document}

\maketitle

\begin{abstract}
% Briefly state the problem, the method, and the key finding.
Gene Regulatory Networks (GRNs) govern critical biological processes like the cell cycle and the p53-Mdm2 DNA damage response. Traditional Boolean modeling of these networks suffers from state-space explosion and lacks differentiability. Here, we introduce an optimized Semi-Tensor Product Graph Neural Network (STP-GNN) framework that provides a continuous, differentiable algebraic state space representation of logic networks. By applying Projected Gradient Descent (PGD) to a 5-node p53 network and a 10-node cell cycle model, we discover that network vulnerabilities are 'holographic'—distributed globally rather than localized in single 'master' edges. These mathematical findings are empirically validated against CRISPR dependency data from the Cancer Dependency Map (DepMap).
\end{abstract}

\section{Introduction}
% Introduce GRNs and the limitations of discrete Boolean simulations.
% Introduce STP and how it solves these limitations.
Gene Regulatory Networks (GRNs) are often modeled as discrete Boolean networks. While conceptually simple, discrete simulations suffer from significant drawbacks: they cannot be optimized using modern gradient-based methods, and they struggle to capture the continuous nature of biological transitions. 

The Semi-Tensor Product (STP) of matrices provides a mathematical bridge, converting logical equations into algebraic equations. However, standard STP operations scale at $\mathcal{O}(4^n)$, causing memory exhaustion even for small networks. This paper introduces an optimized $\mathcal{O}(n \cdot 2^n)$ STP operator and integrates it with deep learning frameworks to perform adversarial vulnerability discovery on the p53 and cell cycle pathways.

\section{Mathematical Foundations}
% Explain STP and ASSR.
% Mention the O(4^n) to O(n * 2^n) optimization trick.
\subsection{Semi-Tensor Product and Algebraic State Space}
The left Semi-Tensor Product, denoted by $\ltimes$, generalizes matrix multiplication. For variables $x \in \{0, 1\}$, we define a mapping to the vector space $\Delta_2$ such that $1 \to [1, 0]^T$ and $0 \to [0, 1]^T$. The global state of an $n$-node network is the Kronecker product of all node states: $x(t) = \ltimes_{i=1}^n x_i(t)$.

\subsection{The Differentiable STP Operator}
The state transition is defined algebraically as $x(t+1) = L \ltimes x(t)$. To make this differentiable, we relax the Boolean logic using a temperature-scaled softmax over the transition parameters $\Theta$. We implement an Implicit Vector-Jacobian Product (VJP) that bypasses the explicit construction of the $2^n \times 2^n$ Jacobian matrix, achieving exponential speedups and enabling analysis of larger networks.

\section{Methods}
% Describe the PGD attack and the DepMap validation setup.
\subsection{Adversarial Network Perturbation}
We subjected the logical transition matrices to Projected Gradient Descent (PGD) to find the minimal structural perturbation ($\epsilon_{critical}$) that forces the network into an arrested state (e.g., G1/S arrest). 

\subsection{Empirical Validation via DepMap}
Predictions derived from the 10-node cell cycle model were cross-referenced with empirical CRISPR knockout data from the Cancer Dependency Map (DepMap) to verify the biological relevance of the mathematically identified bottlenecks.

\section{Results}
% Summarize the p53 "Holographic" finding.
% Summarize the Cell Cycle DepMap validation.
\subsection{The p53 Network and Holographic Vulnerability}
In the 5-node p53-Mdm2 network, the PGD attack discovered a critical collapse threshold at $\epsilon = 2.99$. Strikingly, targeted "knock-in" mutations of the top 5 most vulnerable edges failed to induce apoptosis. The system only collapsed when the perturbation was applied globally, revealing that the network's robustness is a distributed, "holographic" property.

\subsection{Cell Cycle Bottlenecks and Empirical Parity}
In the 10-node cell cycle model, the attack identified the CycD/E2F transition as the primary bottleneck. DepMap validation confirmed this theoretical finding: cancer cell lines with RB1 loss exhibit a statistically significant hyper-dependency on E2F1, proving that the algebraic bottlenecks identified by the STP-GNN govern survival in real-world tumors.

\section{Conclusion}
% Wrap up.
By unifying algebraic logic with gradient descent, the STP-GNN framework overcomes the limitations of discrete simulations. Our findings suggest that effective therapeutic interventions must account for the holographic topology of GRNs, moving beyond single-target paradigms to structural, network-wide disruption strategies.

\end{document}
